% Transcription — fonds Alexandre Grothendieck, shelfmark 15, batch 2 (pages 21-40).
% Established from the facsimile archives/batches/15/batch-02.pdf (a mirror of
% https://grothendieck.umontpellier.fr/15.pdf), per the skill
% transcribe-grothendieck. The batch file carries no cover sheet: PDF page N
% is page 20+N of the folder (the Montpellier footer prints « 21 » ... « 40 »),
% and the archivists' pencil numbers bottom-left agree on every transcribed
% leaf (21, 22, 23, 25, 27, 29, 31, 33, 35, 37, 39 are legible).
%
% Pass: Opus 5.5 (claude-opus-5-5), 2026-09-23 — first pass, unchecked against the pages by a human.
%
% Eleven of the twenty pages are transcribed. Absent: page 24, blank; pages
% 26, 28, 30, 32, 34, 36, 38 and 40, the typed versos of pages 25-39 — leaves
% of an unrelated typescript on holomorphic functions of one complex variable
% (« Chapitre II », Cauchy's theorem, analytic continuation), nothing of it
% part of this folder's mathematics.
%
% What the batch is. Two runs in his hand, neither paginated by him.
%   (1) Pages 21-23, white paper, blue ink, continuing batch 1 (page 21 opens
%   mid-sentence, « gerbe est alors une catégorie A'-linéaire … »):
%   21     canonical invertible objects det(X_{A'}) via a reduced norm
%          N_red : σ* → A*, σ = End(X); multiplicativity of N_red; the
%          conditions Pic(A) = 0, Φ ∩ A* = {1}, Φ·A* = K* (Φ = N_red(E*)),
%          verified for A = Z; a canonical object of degree … of Rep(G_{A'}).
%   22     the band reduced to μ_2; C = category of elliptic curves in char. p
%          (supersingular, one presumes): by Deuring a Deuring category over
%          Z, attached to the quaternion algebra over Q split outside p and ∞;
%          a gerbe with band μ_2 over Z[1/p]; for A' = ∏_{ℓ≠p,∞} Z_ℓ a
%          trivialisation (fibre functor of Rep(Ĝ_{A'})), φ(T) ≃ A'(1); the
%          gerbe H_{A'} of square roots of A'(1).
%   23     G_{A'} ≃ H_{A'}; the invariant: a torsor under H¹(A', μ_2) ≃
%          ∏ Z_ℓ*/Z_ℓ*², then an invariant in Π/ρ(Z/2Z) with
%          ρ : H¹(Z[1/p], μ_2) ≃ Z/2Z × Z/2Z → H¹(A', μ_2), ρ(1,0)_ℓ = 0 or 1
%          according as p (resp. −1) is a square in Z_ℓ*. Two margin notes
%          (no natural choice of origin; eliminating Z/2Z × Z/2Z by
%          « orientations » at ∞ and at p). The page stops half-way down.
%   (2) Pages 25-39, yellowed paper, headed « Cat. de Deuring sur A »:
%   25     definition: A-linear categories C ≃ P(σ/A), three equivalent forms
%          (via an object X_0 with σ = End_C(X_0) and P ↦ P ⊗_σ X_0); σ, σ'
%          give equivalent C iff σ' ≃ End(P), P ∈ P(σ/A); admissible (« de
%          Deuring ») functors F : C → C' over A → A'.
%   27     admissibility tested on one X: pairs (X', u : σ → σ'); the base
%          change C_{A'} as solution of a universal problem,
%          Hom_ad(C, C') ≃ Hom_ad_{A'}(C_{A'}, C'); the Gr-category of
%          autoequivalences = invertible bimodules (« bitorseurs ») under σ.
%   29     equivalences C → C' form a torsor-category; Ex. 1: A Dedekind,
%          E semisimple over K, maximal orders, Deuring categories « de type
%          E »; margin: π_1 = Z* (Z the integral closure of A in the centre
%          of E), = A* if E is central; local case π_0 = 1, π_1 = A*.
%   31     Ex. 2: a norm (ν, α) of C relative to a Deuring category D,
%          ν(P ⊗_σ X) ≃ Nr_{σ/A}(P) ⊗_A ν(X); universal ν_A(C); for fixed X
%          norms ↔ objects T of D; with Nr_{σ/A}(P) ≃ A, an admissible
%          φ : C → C' gives a canonical T_φ with ν(φ(X)) ≃ T_φ.
%   33     Ex. 3: pairs (C, τ), τ ∈ ν_A(C), as generalised gerbes: a gerbe
%          with band μ_n gives (C, τ) (representations of degree n and rank 1,
%          τ the determinant); autoequivalences of (C, τ) = bimodules P with
%          N_{σ/A}(P) ≃ A; in the local case a strictly commutative Picard
%          category with π_1 = μ_n(A), equivalent to Tors(S, μ_n).
%   35     0 → H¹(S, μ_n) → π_0(Brd(C, τ)) → Ram(C/A) → 0,
%          Ram(C/A) := Brd(C)/Pic(A) ≃ Σ_p Z/e_p(C)Z (ramification degrees);
%          equivalences (C_{A'}, T_{A'}) → (C', T') form a torsor-category
%          under Éq(C_{A'}, T_{A'}) ≃ Éq(C', T'), under Tors(S', μ_n) when C'
%          is unramified; margin « À vérifier !! ».
%   37     the Deuring-Brauer category C carries a canonical T ∈ ν(C) (section
%          c of Éq(C, T) → Éq(C)); isomorphism classes of equivalences form a
%          torsor under Brd(C', T')/Im Brd(C); the maps
%          Brd(C) → Brd(C, T) → Brd(C', T'), ρ ↦ (ρ, u_ρ) ↦ (ρ', u_{ρ'}).
%   39     the case of interest (ξ ∈ Br(Q) ramified only at p (?) and ∞):
%          Brd(C) is the Picard category Z/2Z, generated by the Frobenius
%          quotient E ↦ E^{(p)} = E/Ker F_E, i.e. the prime ideal m of σ over
%          p, N(m) = pZ, ρ̄ = p; over A' = Z_p an extension of Z/2Z by
%          Z_p*/Z_p*² × Z/2Z; over A' = R, H¹(R, μ_2) = R*/R*². The page ends
%          mid-sentence, « On trouve donc », continued presumably in batch 3.
%
% « Deuring ». The recurrent name, written « Deu—ring » with a long flourish
% (pp. 22, 25-39) and once « Deuring-Brauer » (p. 37), is read Deuring: page
% 22 says « D'après Deuring » of the category of elliptic curves in char. p
% and its quaternion algebra, which settles it.
%
% Notation fixed for the batch. σ for his « o »-like letter (σ = End(X), an
% A-algebra; his σ' the commutant or the algebra after base change);
% \mathbb{P}(σ/A) for his double-struck P, the category of right σ-modules
% locally free over Spec A as he glosses it on page 25; N_red, Nr_{σ/A},
% N_{σ/A} as he writes them (he alternates); ν, ν_A(C) the norm and its
% universal target; \mathcal{G}_{A'}, \mathcal{H}_{A'} the gerbes; Éq for his
% underlined « Éq »; Brd for the underlined sigle of pages 35-37 (read « Brd »,
% possibly « Bird »; the reading of the sigle is doubtful); Ram, Tors, Rep.
% Plain square brackets in the text are his: either his own brackets, or
% words he added in the interline; editorial additions are \add{}.
%
% Legibility. The statements and formulas carry; the fast connective prose,
% notably on pages 21, 29, 33 and 35, and the struck passages are partly
% \ill{}. Nothing on these leaves is dated.
%
% The dating is the inventory's own for the folder, copied verbatim. Its
% group, [10-18] « Théorie des motifs », is dated [à partir de 1958-à partir
% de 1982].
\documentclass[11pt,a4paper]{article}
% Le préambule commun à toutes les transcriptions du fonds Grothendieck.
%
% Il tient en une idée : **l'incertitude se compose, elle ne se résout pas.**
% Une transcription qui lisse un mot douteux a détruit la seule chose pour
% laquelle on la faisait — un lecteur doit pouvoir savoir, à chaque endroit,
% ce qui était sur le feuillet et ce qui a été ajouté. Les six macros qui
% suivent sont donc l'appareil critique, et non de la décoration.
%
% Les mêmes commandes sont reconnues par `scripts/render.mjs`, qui produit la
% vue de lecture du site. Une macro ajoutée ici doit l'être aussi là-bas,
% faute de quoi le rendu échouera bruyamment — ce qui est voulu : un rendu
% silencieusement faux est pire qu'un rendu absent.

\ProvidesPackage{grothendieck}[2026/08/08 Transcriptions du fonds Grothendieck]

\usepackage{amsmath,amssymb,amsthm}
\usepackage{mathrsfs}
\usepackage{stmaryrd}
% Les diagrammes commutatifs sont partout dans ce fonds : c'est la langue
% même de la géométrie algébrique de Grothendieck, et une transcription qui
% les rendrait en prose ne transcrirait rien.
\usepackage{tikz-cd}
% Français par défaut — c'est la langue des feuillets — mais l'anglais est
% chargé aussi : la traduction et le résumé sont des documents anglais, et la
% typographie française (espaces avant les deux-points) y serait une faute.
% Ces documents basculent par \selectlanguage{english} en tête de corps.
\usepackage[english,french]{babel}
\usepackage{fontspec}
\usepackage{geometry}
\geometry{a4paper,margin=25mm,marginparwidth=28mm}
\usepackage{marginnote}
\usepackage{xcolor}
% ulem, not soul. `soul`'s \st and \ul are built on a character-by-character
% scan that XeLaTeX's font machinery breaks: the document fails with an
% undefined control sequence rather than degrading. ulem does the same two jobs
% and is Unicode-engine safe. `normalem` stops it hijacking \emph, which the
% transcriptions use for Grothendieck's own emphasis.
\usepackage[normalem]{ulem}
\usepackage{hyperref}
\hypersetup{colorlinks=true,linkcolor=black,urlcolor=black,pdfborder={0 0 0}}

% --- Métadonnées du lot -----------------------------------------------------
% Reprises telles quelles par le script de rendu, qui les affiche en tête de
% la vue de lecture. Elles fixent ce que le fichier prétend couvrir : sans
% elles, une transcription flotte sans point d'attache dans un fonds de seize
% mille feuillets.

\newcommand{\folder}[1]{\gdef\grothFolder{#1}}
\newcommand{\batch}[1]{\gdef\grothBatch{#1}}
\newcommand{\leaves}[2]{\gdef\grothFirst{#1}\gdef\grothLast{#2}}
\newcommand{\foldertitle}[1]{\gdef\grothFoldertitle{#1}}
\gdef\grothFolder{?}\gdef\grothBatch{?}\gdef\grothFirst{?}\gdef\grothLast{?}\gdef\grothFoldertitle{}

% --- Le résumé --------------------------------------------------------------
%
% Il ouvre la lecture modernisée, sous le titre « Résumé », et il est composé
% en petit. Ce n'est pas un ornement : c'est le seul passage du document qui
% ne soit pas des mathématiques, et un lecteur doit pouvoir le reconnaître
% avant de l'avoir lu — puis le sauter s'il connaît déjà le sujet, ou s'y
% arrêter s'il ne le connaît pas. Le filet qui le suit marque la reprise du
% corps.

\newenvironment{resume}
  {\small\setlength{\parskip}{0.45em}}
  {\par\medskip\hrule\medskip}

% --- Mots-clés ---------------------------------------------------------------
%
% En anglais, en clôture du résumé de la lecture modernisée : le vocabulaire
% moderne sous lequel chercher ce que les feuillets construisent. Le manifeste
% du site (`npm run manifest`) les extrait pour étiqueter la cote entière —
% cette ligne est la source unique des tags, il n'y a pas de fichier à côté.
\newcommand{\keywords}[1]{\par\smallskip\noindent
  {\footnotesize\textsc{Keywords} --- \textit{#1}}\par}

% --- L'appareil critique ----------------------------------------------------

% Le numéro de feuillet, en marge. C'est l'ancre qui permet de régler un
% désaccord sur une formule en regardant *un* feuillet plutôt qu'une chemise
% de six cents. Le site s'en sert aussi pour tourner les pages du fac-similé
% quand on fait défiler la transcription.
\newcommand{\leaf}[1]{%
  \marginnote{\footnotesize\color{gray!70}\textsf{#1}}%
  \label{leaf:#1}\ignorespaces}

% Illisible : on ne devine pas. Un mot inventé qui se lit comme les autres est
% le pire résultat possible.
\newcommand{\ill}{\textcolor{red!60!black}{[\,\ldots\,]}}

% Lecture douteuse : le mot est donné, et signalé comme incertain.
\newcommand{\uncertain}[1]{\uline{#1}}

% Ajout de l'éditeur — une lettre manquante, un mot sous-entendu.
\newcommand{\add}[1]{\textcolor{blue!50!black}{[#1]}}

% Biffé par Grothendieck lui-même. Ce qu'il a rayé fait partie du document :
% une rature dit souvent où la pensée a changé de direction.
\newcommand{\struck}[1]{\sout{#1}}

% Note du transcripteur, sur l'état matériel du feuillet.
\newcommand{\note}[1]{\par\smallskip{\small\color{gray!60!black}\textit{[#1]}}\par\smallskip}

% Note marginale de Grothendieck, distincte du corps du texte.
\newcommand{\marginal}[1]{\par\smallskip\noindent\hspace{1em}%
  {\small\color{gray!60!black}#1}\par\smallskip}

% --- Titre ------------------------------------------------------------------

\AtBeginDocument{%
  \begin{center}
    {\large\bfseries Fonds Alexandre Grothendieck --- cote n\textsuperscript{o}~\grothFolder}\\[2pt]
    {\itshape\grothFoldertitle}\\[4pt]
    {\small feuillets \grothFirst--\grothLast\ (lot \grothBatch)}\\[2pt]
    {\footnotesize\color{gray!60!black}Transcription établie d'après le fac-similé numérisé
     par l'Université de Montpellier.\\
     Les lectures incertaines sont soulignées, les illisibles notées [\,\ldots\,],
     les ajouts entre crochets.}
  \end{center}
  \vspace{1em}}

\endinput


\folder{15}
\batch{2}
\pages{21}{40}
\dating{1965-[vers 1977]}
\watermark{Édition de démonstration}
\foldertitle{Théorie arithmétique. Théorie de Galois des motifs (1965) : notes manuscites (s.d.), tapuscrit (s.d.).}

\begin{document}

\page{21}
\note{le feuillet continue une phrase commencée au feuillet 20 (lot 1)}

gerbe est alors une [catégorie] $A'$-linéaire, \ill{}. \struck{Soit $X \in$}
Nous dirons,
\struck{Supposons que pour \ill{} Soit $X \in \operatorname{Ob} C$, et
considérons $\sigma = \operatorname{End}(X)$, et}
\struck{$N_{\mathrm{red}} : \sigma^{*} \to A^{*}$}
\note{les trois lignes biffées sont en outre encadrées d'un trait qui les
isole du reste}
sous certaines conditions, on définira les objets inversibles
canoniques, en établissant un système transitif d'isomorphismes entre les
$\det(X_{A'})$ $(X \in \operatorname{Ob} C)$. Pour ceci il suffit de définir
un isom.\ canonique
\[
  \det P \simeq A \qquad (\alpha)
\]
pour les $A$-modules \ill{} à droite $P$ sous un $\sigma =
\operatorname{End}(X)$ $(X \in \operatorname{Ob} C)$, fonctoriel en $P$ (ce
qui \uncertain{suppose} que $N_{\mathrm{red}} : \sigma^{*} \to A^{*}$ est à
valeurs dans $1$, \ill{} $\sigma'$ étant le commutant de $\sigma$ dans $P$
…) et multiplicatif \ill{}
\uncertain{($Q$ sous $\sigma'', \sigma'$, $P$ sous $\sigma', \sigma$)}
\[
  N_{\mathrm{red}}(Q \otimes_{\sigma'} P, \sigma) \simeq
  N_{\mathrm{red}}(Q, \sigma') \otimes N_{\mathrm{red}}(P, \sigma).
\]
\ill{} doit être compatible avec les isom.\ \ill{}. Cela \uncertain{marche}
si on suppose [Prenons $\Phi = N_{\mathrm{red}}(E^{*}) \subset K^{*}$, où
$\sigma = \operatorname{End}(X)$, $E = \sigma \otimes_{A} K$]
\[
  \begin{cases}
    \operatorname{Pic}(A) = 0 \\
    \Phi \cap A^{*} = \{1\}, \quad \Phi \cdot A^{*} = K^{*}
  \end{cases}
\]
\note{sur la seconde ligne de l'accolade, une première version biffée :
\struck{$N_{\mathrm{red}}(E^{*}) \cap N_{\mathrm{red}}(\sigma)$}}

Ces conditions sont vérifiées \struck{par ex.\ si $A = \mathbb{Z}$, \ill{}}
si $A = \mathbb{Z}$, $E_{K}$ $(= E_{\mathbb{R}})$ \ill{}
\struck{\ill{} \ill{} \ill{}} sur $\mathbb{R}$. Donc dans ce cas, on a un
objet inv.\ canonique de degré \ill{} de $\operatorname{Rep}(\mathcal{G}_{A'})$,
de \ill{} jusqu'à que le lien \uncertain{structural} de la gerbe
$\mathcal{G}_{A'}$

\page{22}
se trouve de façon naturelle réduit à $\mu_{2}$.

On peut maintenant \ill{}
\[
  C = \text{catégorie des courbes elliptiques de}
\]
\ill{} (Hom \ill{} sur \uncertain{corps alg.\ clos} de car.~$p$).

D'après Deuring, c'est une catégorie de Deuring [définie sur $\mathbb{Z}$],
\struck{dont} correspond à la classe d'isomorphie des \uncertain{algèbres}
quaternions de degré $4$ (corps de quaternions) sur $\mathbb{Q}$ dont
déployés en toutes les places sauf $p$ et $\infty$ (il est donc d'ordre $2$
dans $\operatorname{Br}(\mathbb{Q})$).

Les torseurs donc, comme \ill{} en termes de \struck{\ill{}} $C$ \ill{}
[la donnée d'un tel corps sur $\mathbb{Q}$, une
gerbe sur $\operatorname{Spec} \mathbb{Z}[\tfrac{1}{p}]$ de
\struck{groupe} lien $\mu_{2}$ [$\simeq$ \ill{} $\times$ $\widehat{\mathcal{G}}$
de inversibles] + trivialisation de ce déterminant, de rang $2$].
\note{« la donnée » est récrit dans l'interligne au-dessus d'une première
version biffée ; « inversibles » est ajouté sous la ligne ; la lecture du
crochet intérieur est incertaine}
On trouve de plus, si
\[
  A' = \prod_{\ell \neq p, \infty} \mathbb{Z}_{\ell},
\]
une banalisation de $C_{A'}$, i.e.~
\struck{représentation de $\mathcal{G}_{A'}$ de}
i.e.\ une trivialisation de $\widehat{\mathcal{G}}_{A'}$,
\struck{soit $E$, (avec un isom.\ \ill{} de $\mathcal{G}_{A'}$ de rang $2$,
degré $1$), $\det E \simeq T$}
[i.e.\ foncteur fibre de $\operatorname{Rep}(\widehat{\mathcal{G}}_{A'})$ sur
$A'$], d'où un isom.
\[
  \varphi(T) \simeq A'(1) \Bigl(\overset{\mathrm{df}}{=}
  \prod_{\ell \neq p, \infty} \mathbb{Z}_{\ell}(1)\Bigr).
\]
\note{une première version biffée précède la formule : \struck{$\varphi(E)$}}

\struck{\ill{} Cela fait donc} En d'autres termes, si
$\mathcal{H}_{A'}$ est la gerbe des racines carrées de $A'(1)$, qui est une
gerbe sous $\mu_{2}$, on a une équivalence canonique de

\page{23}
\[
  \mathcal{G}_{A'} \simeq \mathcal{H}_{A'}
\]
Comme invariant ($C$ dans $\mathcal{G}_{A^{*}}$ [de $\mathcal{G}_{A'}$]
\ill{} ici précisé) on trouve par cette donnée [à isom.\ près]
\struck{explicitement} un torseur sous
\[
  \mathrm{H}^{1}(A', \mu_{2}) \simeq \prod_{\ell \neq p, \infty}
  \mathbb{Z}_{\ell}^{*}/\mathbb{Z}_{\ell}^{*2},
  \qquad
  \mathbb{Z}_{\ell}^{*}/\mathbb{Z}_{\ell}^{*2} \simeq
  \begin{cases}
    \mathbb{Z}/2\mathbb{Z} & \text{si } \ell \neq 2 \\
    \mathbb{Z}/2\mathbb{Z} \times \mathbb{Z}/2\mathbb{Z} & \text{si } \ell = 2
  \end{cases}
\]
\note{avant la formule, une première version biffée :
\struck{$\mathrm{H}^{1}(\mathbb{Q}$ \ill{} $A', \mu_{2}) \simeq$ \ill{}}}

\marginal{NB il n'y a absolument pas de façon naturelle de choisir une
origine \struck{\ill{}} « \uncertain{a priori} » \struck{\ill{} dans}
\struck{\ill{} dans} \ill{} dans $\mu_{\ell}$ mod
$(\mathbb{Z}/\ell\mathbb{Z})^{*2}$ $(\ell \neq 2)$, $\mu_{8}$ mod
$(\mathbb{Z}/8\mathbb{Z})^{*2}$}
\note{note marginale écrite de biais dans la marge gauche, en partie
encadrée ; l'exposant porté sur $\mu_{\ell}$ n'est pas lu}

Si on tient compte du fait que la gerbe $\mathcal{G}_{A'}$ \ill{}
déterminée seulement à équivalence (non unique près), \ill{} la \ill{}
d'une classe dans l'isomorphie d'équivalences \ill{}, \ill{} donc d'un
élément de
\[
  \mathrm{H}^{1}\bigl(\underbrace{\mathbb{Z}[\tfrac{1}{p}]}_{A}, \mu_{2}\bigr)
  \simeq \mathbb{Z}/2\mathbb{Z} \times \mathbb{Z}/2\mathbb{Z}
\]
(l'élément $(1,0)$ \struck{\ill{} de $\mathbb{Z}/2\mathbb{Z}$} [resp.~
$(0,1)$] correspondant resp.\ à $\sqrt{p}$ [resp.\ $\sqrt{-1}$]) on trouve un
invariant \struck{dans}
\marginal{on peut éliminer le gpe $\mathbb{Z}/2\mathbb{Z} \times
\mathbb{Z}/2\mathbb{Z}$ grâce à la donnée de la structure
\uncertain{s-pl.} : a) orientation à l'$\infty$, b) « orientation » en $p$}
\[
  \Pi/\rho(\mathbb{Z}/2\mathbb{Z}),
\]
où $\Pi$ est le torseur sous
$\mathrm{H}^{1}(A', \mu_{2}) = \prod_{\ell \neq p, \infty}
\mathbb{Z}_{\ell}^{*}/\mathbb{Z}_{\ell}^{*2}$ des trivialisations de
$\mathcal{H}_{A'}$, et
\[
  \rho : \mathrm{H}^{1}(\mathbb{Z}[\tfrac{1}{p}], \mu_{2}) \simeq
  (\mathbb{Z}/2\mathbb{Z}) \times \mathbb{Z}/2\mathbb{Z} \hookrightarrow
  \mathrm{H}^{1}(A', \mu_{2})
\]
l'homom.\ naturel, défini par la condition [(resp.\ $\rho(0,1)_{\ell}$)]
\[
  \rho(1,0)_{\ell} =
  \begin{cases}
    0 & \text{si } p \text{ [resp.\ } -1\text{] est un carré dans }
        \mathbb{Z}_{\ell}^{*} \\
    1 & \text{sinon}
  \end{cases}
\]
($\Longleftrightarrow$ $p$ carré dans $(\mathbb{Z}/\ell\mathbb{Z})^{*}$ si
$\ell \neq 2$).
\note{le « $\ell \neq 2$ » est d'abord écrit sous la parenthèse, biffé, puis
récrit au-dessus ; le feuillet s'arrête à mi-hauteur}

\section{Cat. de Deuring sur A}
\note{titre de sa main, souligné, en tête du feuillet 25 ; la suite (feuillets
25 à 39, rectos d'un papier jauni) n'est pas paginée par lui. Le nom, tracé
« Deu—ring » avec un long paraphe, est lu Deuring d'après le feuillet 22}

\page{25}
C'est une catégorie prédditive \uncertain{$A$-linéaire} telles que
\[
  \exists\ A\text{-algèbre } \sigma, \text{ avec } C \approx
  \mathbb{P}(\sigma/A)
\]
\note{« $A$- » est récrit sur un premier mot biffé devant « algèbre »}
[$\mathbb{P}(\sigma/A)$ : modules $P$ à droite sur $\sigma$ loc.\ libres de
rang $1$ \uncertain{sur} $\operatorname{Spec} A$] [équiv.\ $A$-linéaire]
\[
  \Updownarrow
\]
$\exists\, X_{0} \in \operatorname{Ob} C$ t.q.\ a) pour tt
$P \in \operatorname{Ob} \mathbb{P}(\sigma/A)$, où $\sigma =
\operatorname{End}_{C}(X_{0})$, $P \otimes_{\sigma} X_{0}$ existe ;
b) $P \mapsto P \otimes_{\sigma} X_{0}$ une équivalence de catégories
\[
  \Updownarrow
\]
$\forall\, X_{0} \in \operatorname{Ob} C$, a) b).

\emph{NB} À équiv.\ (canon.\ unique) près, $C$ est donc défini par une
algèbre $\sigma$. Pour que $\sigma, \sigma'$ \struck{définissent} [définissent]
des $C$ équivalentes, il faut et suffit qu'il existe
$P \in \mathbb{P}(\sigma/A)$, avec $\sigma' \simeq \operatorname{End}(P)$,
i.e.\ que $\sigma'$ soit obtenu en « tordant par autom.\ intérieurs par un
$P \in \mathbb{P}(\sigma/A)$ ». Si $\sigma$ est commutatif (\ill{} \ill{}
\ill{} !) on trouve que $\sigma \simeq \sigma'$.

\struck{\ill{}}, Soient $C, C'$ deux catégories [$A$-lin.\ avec
$A \simeq A'$ donné] de Deuring. Un foncteur
$F : C \to C'$ \struck{\ill{}} [$A$-lin.] est dit \emph{admissible}
[(\uncertain{ou de Deuring})] s'il \uncertain{prolonge} commutant à
l'opération $P \otimes_{\sigma} X$ pour $X \in \operatorname{Ob} C$,
$\sigma \in \operatorname{End}_{C}(X)$, $P \in \mathbb{P}(\sigma/A)$. Je
\uncertain{dis} \ill{}
\note{l'interligne au-dessus de « deux catégories » porte « $A$-lin.\ avec
$A \to A'$ donné », dont la flèche est écrite comme un $\simeq$ ; lecture
incertaine}

\page{27}
de l'exiger pour un $X$, et pour $X$ fixé, la \struck{\ill{} de}
catégorie des foncteurs est équivalente à la catégorie des couples
$(X', u)$ où $X'$ est un objet de $C'$,
\[
  u : \sigma = \operatorname{End}(X) \longrightarrow \sigma' =
  \operatorname{End}(X')
\]
un homom.\ compatible avec $A \to A'$, tel que
$\forall\, P \in \operatorname{Ob} \mathbb{P}(\sigma/A)$,
\struck{$P \otimes$} posant $P' = P \otimes_{\sigma} \sigma' \in
\mathbb{P}(\sigma'/A')$, le produit tensoriel $P' \otimes_{\sigma'} X'$
existe.

[\emph{NB} le « tel que » est superflu si p.ex.\ $C'$ est de Deuring.]

\struck{Considérons}
Catégorie $C_{A'}$ ($C$ catégorie de Deuring sur $A$ ; correspondant à
solution d'un pb universel)
\[
  \Bigl[\ \operatorname{Hom}_{\mathrm{ad}_{A \to A'}}(C, C') \simeq
  \operatorname{Hom}_{\mathrm{ad}_{A'}}(C_{A'}, C')\ \Bigr]
\]

\struck{Si $C, C'$ sont des catégories de Deuring \ill{} deux, pour que}
Si $C_{A'} \to C'$ sont une équivalence, il faut que $C'$ soit de Deuring
aussi, et que $\sigma \otimes_{A} A' \to \sigma'$ soit un isomorphisme.

La Gr-catégorie des autoéquivalences de $C$ [$\approx$] est la catégorie
des bitorseurs $P$ sous $\sigma, \sigma$, i.e.\ \uncertain{loc.\ sur $A$
iso.}\ pour les deux structures, avec le produit tensoriel évident.
\struck{Si} $A = A'$

\page{29}
la catégorie des équivalences de $C$ avec $C'$ ($C'$ étant de Deuring) est
un Tors-catégorie sous la Gr-catégorie précédente.

\emph{Ex 1} Soit $A$ anneau de Dedekind, [corps des fractions $K$],
\struck{\ill{} $A$ \ill{}} donc l'isomorphie d'algèbres
[\uncertain{pures}] abs.\ semi-simples $E$ sur $K$, on \ill{} une
\ill{} d'équivalences de catégories de Deuring : les catégories
[de Deuring de type $E$]
\[
  \mathbb{P}(\sigma'/A) \approx \mathbb{P}(\sigma/\text{ordre max.\ de }
  \sigma)
\]
définies par les ordres \emph{maximaux} de $E$ : \struck{\ill{}}
\ill{} extension de base $A \to A'$, où $A'$ un corps ou un anneau de
Dedekind, cat.\ résiduelles triviales \ill{}, complété \ill{} ! !). Si
$A$ est local, $E$ \struck{central} \ill{} [\ill{}] de
\marginal{La catégorie des Autoéquivalences d'une telle $C$ a comme
$\pi_{1}$ le groupe $Z^{*}$, où $Z$ est la clôture intégrale de $A$ dans le
cent.\ de $E$. (Donc si $E$ central, $\pi_{1} = A^{*}$)}
\struck{\ill{}} [$\mathrm{Gr}$-] alors la catégorie \struck{\ill{}} des
autoéquivalences d'une $C$ de Deuring de type $E$ est un $\pi_{0} = 1$
(les bimodules sur $\sigma$ inversibles des deux côtés sont triviaux).
\struck{\ill{} la catégorie} [Mais] il vient $\pi_{1} = A^{*}$.
\note{« Mais » est ajouté au-dessus du passage biffé ; la suite des
feuillets 29 à 39 est écrite dans la même encre brune, avec des reprises à
l'encre noire}

\page{31}
\emph{Ex.\ 2} Soit $C$ comme dans Ex.\ 1 et soit $D$ une catégorie de
Deuring \ill{} sur $A$. Une \emph{norme} de $C$ relative à $D$ est un couple
$(\nu, \alpha)$, où
\[
  \nu : C \longrightarrow D
\]
un foncteur (\ill{} \ill{} admissible !),
\[
  \alpha = (\alpha_{X, P}) : \nu(P \otimes_{\sigma} X)
  \xrightarrow{\ \sim\ } \mathrm{Nr}_{\sigma/A}(P) \otimes_{A} \nu(X)
\]
un système d'isomorphismes, satisfaisant [non évidentes] à des conditions de
transitivité, et fonctorialité en $P, X$ … .
\note{la lettre notée ici $D$ est d'abord tracée comme un $\mathcal{P}$ ;
la suite ($\nu : C \to D$) la fixe}

\marginal{Généralisation au cas où $D$ n'est plus de Deuring, mais seulement
$A$-lin.\ et tq les $\mathrm{N}_{\sigma/A}(P) \otimes_{A} Y$
$(Y \in \operatorname{Ob} D)$ existent. — Montrer qu'il y a une $D$
universelle recevant une $(\nu, \alpha)$ universelle, notée $\nu_{A}(C)$.
\struck{\ill{}} Noter que $\nu_{A}(C)$ \ill{} aux conditions \ill{} des
\ill{} dans le cas \ill{} de Ex.\ 1}
\note{note marginale écrite de biais, à l'encre noire, dans la marge gauche
du haut du feuillet ; un trait double la sépare de sa seconde moitié}

\struck{\ill{}} Pour $X$ fixé dans $\operatorname{Ob} C$, la cat.\ des
normes \struck{\ill{}} est équivalente à celle des objets $T_{0}$ de $D$,
en définissant $\nu$ par
\[
  \nu(P \otimes_{\sigma} X) \simeq \mathrm{Nr}_{\sigma/A}(P) \otimes T.
\]
\struck{$(\nu, \alpha)$} Supposons maintenant qu'il existe un isom.~
fonctoriel en $P \in \mathbb{P}(\sigma/A)$
\[
  \mathrm{Nr}_{\sigma/A}(P) \simeq A
\]
\marginal{avec conditions de transitivité convenables}
\note{sous la formule, une ligne biffée : \struck{$A' \to A''$}}

\marginal{Notons que la catégorie $\nu_{A}(C)$ est munie d'une section
canonique $T$ et $\nu_{C}$ \ill{} \ill{}}

Soit $A' \longleftarrow A$ un chgt de base comme dans ex.\ 1, et considérons
sur $A'$ des catégories de Deuring $C', D'$, et une norme
$(\nu, \alpha) : C' \to D'$. Soit $\varphi : C \to C'$ un foncteur
admissible. Alors il y a un objet can.\ $T_{\varphi}$ de $D'$, et un isom.
\[
  \nu(\varphi(X)) \simeq T_{\varphi}.
\]

\begin{tikzcd}
C \arrow[r, "\varphi"] & C' \arrow[r, "\nu"] & D' \\
A & A' &
\end{tikzcd}

\note{sous le diagramme, $A$ et $A'$ sont écrits sous $C$ et $C'$, sans
flèche ; l'interligne porte « et \ill{} » au-dessus de $\varphi$ et « un »
au-dessus de $\nu$}

\page{33}
\emph{Ex.\ 3} Soit, en plus de la catégorie de Deuring $C$ sur $A$ comme
dans Ex.\ 2, un \struck{$C$} élément $\tau$ de $\nu_{A}(C)$. Si $E$ est
[\uncertain{central}] de rang $n$, il faut considérer que les $C$
généralisent les \struck{(\ill{} \ill{} \ill{} finis)} gerbes sur $A$
trivialisées par \ill{}, \struck{\ill{} dans \ill{}, \ill{} \ill{} de
\ill{}} \ill{} $E$ \ill{} \ill{} [\ill{} \ill{}] de $\operatorname{Br}(K)$,
i.e.\ le classe de $C$ dans \ill{} \ill{} \ill{} et que les couples
$(C, \tau)$ \ill{} la donnée d'une gerbe de groupe $\mu_{n}$.

\marginal{gerbe sur $A$ de lien $G_{m}$ $\mapsto$ catégorie $C$ des
représentations de celle-ci \ill{}, poids $1$. $C \mapsto$ gerbe des
foncteurs admissibles de $C$ dans les modules \ill{} (trivialisations de
$C$) \ \ \ cat.\ de Deuring type $E$ \ill{} $\xi \in$
$\operatorname{Im}(\operatorname{Br} A \to \operatorname{Br} K)$}
\note{note marginale de biais, dans la marge gauche, reliée par une
accolade au passage précédent}

Pour ceci, \struck{\ill{}} notons \struck{la catégorie des \ill{}
équivalentes} que \struck{pour} une gerbe \struck{des} liens $\mu_{n}$ sur
$A$ définit bien une $(C, \tau)$ :
$C =$ les \struck{banalisations de} représentations de degré $n$, de rang
$1$ de la gerbe, $\tau =$ \struck{représentation déterminant}.
Nous allons déterminer d'autre part les autoéquivalences de $(C, \tau)$ ; on
se donne \ill{} d'un objet $X$ de $C$, [d'où $\sigma =
\operatorname{End}(X)$, i.e.\ \ill{}] la catégorie des \struck{\ill{}}
bimodules $P$ sous $\sigma$, inv.\ loc.\ sur $A$ des deux côtés,
\uncertain{munis} d'un isom.
\struck{$\nu(X) = L \otimes$}
\[
  \mathrm{N}_{\sigma/A}(P_{\sigma}) \simeq A
  \qquad (\text{i.e.\ } \tau = L \otimes_{A} \nu(X))
\]
\marginal{$\varphi(X) = P \otimes_{\sigma} X$,
$\nu(\varphi(X)) = \mathrm{N}(P) \otimes \nu(X) \simeq \varphi(\nu(X))$}

\struck{compatible avec}
\struck{Supposons que tous les $P$ \ill{} soient isomorphes ($A$ \ill{}
local, sans \ill{}) \ill{} de $\mathbb{P}(\sigma/A)$}
\note{passage encadré et biffé de deux traits obliques ; au-dessus, « de
$\mathbb{P}(\sigma/A)$ » ajouté}
Alors l'\ill{} est une Gr-catégorie, i.e.\ \ill{}
[et à une catégorie de Picard \uncertain{strictement} commutative],
d'objet $1$ dont les éléments $\lambda \in A^{*}$ tels que
$\lambda^{n} = 1$ ; \ill{} de Picard commutative \ill{}
$\mathrm{Tors}(S, \mu_{n})$ des

\page{35}
torseurs \struck{\ill{}} sur $S$ de groupe $\mu_{n}$, et \uncertain{son
$\pi_{0}$} \struck{\ill{}} s'insère dans
\[
  0 \longrightarrow \mathrm{H}^{1}(S, \mu_{n}) \longrightarrow
  \pi_{0}\bigl(\mathrm{Brd}(C, \tau)\bigr) \longrightarrow
  \mathrm{Ram}(C/A) \longrightarrow 0
\]
\note{le premier terme est surchargé au point d'être illisible ; le dernier,
d'abord un mot biffé, est remplacé par « $\mathrm{Ram}(C/A)$ » dans
l'interligne. Le sigle « $\mathrm{Brd}$ », souligné, est lu ainsi aux
feuillets 35 et 37 ; la lecture en est douteuse}

(\struck{Si $\operatorname{Pic}(A) = 0$, \ill{} la catégorie}
où
\[
  \mathrm{Ram}(C/A) \overset{\mathrm{df}}{=}
  \mathrm{Brd}(C)/\operatorname{Pic}(A) \simeq \sum_{p}
  \mathbb{Z}/e_{p}(C)\mathbb{Z}.
\]
\marginal{$e_{p}$ : degré de ramification}

\struck{Mais} Soit maintenant $A'$ une $A$-algèbre (commutative), [et
$T' \in \nu_{A'}(C')$],
\struck{\ill{} \ill{} \ill{}}
et $C'$ une catégorie de Deuring sur $A'$, supposons
$(C_{A'}, T_{A'}) \approx (C', T')$.
\struck{Alors l'ensemble des}
\note{la parenthèse porte d'abord « $(C_{A'}, T \approx (C', A')$ », peut-être
une faute de plume}

Il suffit pour \uncertain{cela} que $C_{A'} \simeq C'$ et
$\operatorname{Pic}(A') = 0$, (\ill{} p.ex.\ $A'$ local) Alors l'\ill{} des
équivalences
\[
  (C_{A'}, T_{A'}) \dashrightarrow (C', T'),
\]
i.e.\ des couples $(\varphi, u)$ avec $\varphi$ \struck{équivalence}
\[
  C_{A'} \xrightarrow{\ \varphi\ } C' \quad A\text{-lin.\ admis.}
\]
\struck{\ill{}}
$\varphi : C_{A'} \xrightarrow{\sim} C'$, et
$\nu(\varphi)(T_{A'}) \xrightarrow{\sim} T'$, est une catégorie
\emph{torseur} sous la [Gr-]catégorie
\[
  \text{Éq}(C_{A'}, T_{A'}) \ (\text{à droite}) \ \approx\quad
  \text{Éq}(C', T') \ (\text{à gauche}).
\]
Si p.ex.\ $\mathrm{Ram}(C'/A') = 0$ ($C'$ non ramifié sur $A'$)
\struck{\ill{} \ill{}} c'est donc une catégorie-torseur sous
$\mathrm{Tors}(S', \mu_{n})$. En fait, \struck{\ill{}} [c'est]
une équivalence \struck{\ill{}} canonique
\[
  \text{Éq}(C_{A'}, T_{A'}) \approx
  \text{Éq}(C', T')
\]
(comme \ill{} catégories \emph{de Picard}) et par suite
\struck{\ill{}} une structure de torseur [à] tous catégories
\uncertain{près} s'identifie \ill{} …
\marginal{À vérifier !!}
\note{la note marginale est tracée à gauche, entre deux traits verticaux
qui encadrent les quatre dernières lignes}

\page{37}
Signalons qu'il y a par ailleurs canoniquement, \ill{} \emph{canoniques},
\ill{} la catégorie de Deuring-Brauer $C$ (en termes des équivalences de
$C$), un objet $T$ de $\nu(C)$. Cela signifie en particulier que nous avons
relevé le foncteur canonique [de \uncertain{Éq}-cat.]
\[
  \text{Éq}(C, T) \longrightarrow
  \text{Éq}(C)
\]
par une section $c$ (flèche courbe de retour, marquée $c$).
\note{la section est dessinée comme un arc fléché qui revient de
$\text{Éq}(C)$ vers $\text{Éq}(C, T)$,
avec l'étiquette $c$}
\ill{} pour la catégorie $C'$, équivalente à $C$, on obtient un objet
$T_{C'}$ canonique dans $\nu(C')$.

On voit donc que pour $(A', C', T')$ comme plus haut,
\struck{l'ensemble des} tel que $(C', T')$ soit équivalente à
$(C, T)_{A'}$, les équivalences [\struck{catégorie-torseur sous}]
\[
  (C, T)_{A'} \xrightarrow{\ \sim\ } (C', T')
\]
forment une \struck{\ill{}} catégorie-torseur sous
\[
  \text{Éq}\bigl((C, T)_{A'}\bigr) \simeq
  \text{Éq}(C', T')
\]
\struck{\ill{} une catégorie torseur \ill{}}
\struck{$\text{Éq}(C', T') / \text{Éq}(C)$
\uncertain{par} \ill{}}
[les classes d'isom.\ d'objets forment] un torseur sous
\[
  \mathrm{Brd}(C', T') / \operatorname{Im} \mathrm{Brd}(C).
\]
Le torseur ne dépend \ill{} du choix d'un $C$ dans la classe de catégories
\ill{} … . Je \ill{} donc explicite l'application canonique \ill{} :
\[
  \begin{array}{ccccc}
    \mathrm{Brd}(C) & \xrightarrow{\ c\ } & \mathrm{Brd}(C, T) &
      \longrightarrow & \mathrm{Brd}(C', T') \\
    \rho & \longmapsto & (\rho, u_{\rho}) & \longmapsto &
      (\rho', u_{\rho'}) \\
    \wr\!\wr \\
    P_{0} & \longmapsto & \mathrm{N}_{\sigma/A}(P_{0}) \otimes u_{\rho} &
      \longmapsto & \mathrm{N}_{\sigma'/A'}(P'_{0}) \otimes (u_{\rho})'
  \end{array}
\]
\note{le signe entre $\rho$ et $P_{0}$ est un $\simeq$ vertical ; le produit
de la dernière ligne est surchargé}

\page{39}
Dans le cas qui nous intéresse ($\xi \in \operatorname{Br}(\mathbb{Q})$
d'\uncertain{ordre} \ill{}, ramifié seulement en \uncertain{$p$} et $\infty$)
\struck{\ill{} \ill{}} on trouve que $\mathrm{Brd}(C)$ est la
\struck{gpe} [catégorie] de Picard $\mathbb{Z}/2\mathbb{Z}$,
$\mathbb{Z}/2\mathbb{Z}$, le générateur \uncertain{explicité}
d'invariant \ill{} de $\pi_{0}$ \uncertain{lequel} correspond au
\struck{l'automorphisme} foncteur sur les courbes elliptiques
\[
  E \longmapsto E^{(p)} = E/\operatorname{Ker} F_{E}
\]
correspond à l'idéal \uncertain{premier} $\mathfrak{m}$ de $\sigma$ en $p$,
et on a
\[
  \mathrm{N}_{\sigma/\mathbb{Z}}(\mathfrak{m}) \longrightarrow
  \mathrm{N}_{\sigma/\mathbb{Z}}(\sigma) = \mathbb{Z}, \qquad
  \mathrm{N}_{\sigma/\mathbb{Z}}(\mathfrak{m}) = p\mathbb{Z},
\]
on choisit $\bar{\rho} = p$.
\note{le signe lu $p$ dans la parenthèse initiale ressemble à un « 2 » ; le
feuillet 22 place les ramifications en $p$ et $\infty$}

Ceci nous dit que si $p$ est \emph{inversible dans $A'$}, dans l'image de
$(\rho, \bar{\rho})$ dans
\[
  \mathrm{Brd}(C_{A'}, T_{A'}) \simeq \mathrm{H}^{1}(S', \mu_{2})
\]
\struck{\ill{}} est le résidu de $p$ modulo les carrés.

Par contre, prenons le cas où $A' = \mathbb{Z}_{p}$, dans
\struck{\ill{} dans un gpe}
\[
  \mathrm{Brd}(C_{p}, T_{p}) \simeq \mathrm{Brd}(C_{A'}, T_{A'})
\]
est une extension de $\mathbb{Z}/2\mathbb{Z}$ par
$\mathbb{Z}_{p}^{*}/\mathbb{Z}_{p}^{*2} \times \mathbb{Z}/2\mathbb{Z}$
[$= \mathrm{H}^{1}(\mathbb{Z}_{p}, \mu_{2})$], engendré par $(\mathfrak{m}, p)$,
et le \ill{}
$\mathbb{Z}_{p}^{*}/\mathbb{Z}_{p}^{*2}$ $(\simeq \mathbb{Z}/2\mathbb{Z}$
si $p \neq 2)$, on \uncertain{fait} \ill{}
$\mathrm{H}^{1}(\mathbb{Z}_{p}, \mu_{2})$ \struck{\ill{}}.
\marginal{$C_{p} = C'$, $T_{p} = T'$}
\note{le terme $\mathbb{Z}_{p}^{*}/\mathbb{Z}_{p}^{*2}$ de l'extension est
d'abord écrit sous la forme $\mathbb{Z}_{p}^{*}/\mathbb{Z}_{p}^{*2}
(\simeq \mathbb{Z}/2\mathbb{Z}$ si $p \neq 2)$, et
$\mathrm{H}^{1}(\mathbb{Z}_{p}, \mu_{2})$ surcharge un premier terme biffé}

Ceci donne d'innombrables possibilités de
\struck{\ill{}} [\ill{}] trivialiser un torseur sous un
\struck{groupe \ill{}} $\simeq \mathbb{Z}/2\mathbb{Z} \times
\mathbb{Z}/2\mathbb{Z}$ (si $p \neq 2$).

D'autre part, si $A' = \mathbb{R}$, on peut considérer les classes
d'isomorphie de la gerbe de liens $\mu_{2}$ sur $\mathbb{R}$ définie par $C$
(soit $\mathcal{G}_{\mathbb{R}}$, où $\mathcal{G}$ est la gerbe sur
$\mathbb{Z}[\tfrac{1}{p}]$) et l'ensemble des classes d'équivalence avec la
gerbe triviale sur $\mathbb{R}$ ($=$ ensemble des \uncertain{polarisations}
de la gerbe), qui est un torseur sous
\[
  \mathrm{H}^{1}(\mathbb{R}, \mu_{2}) = \mathbb{R}^{*}/\mathbb{R}^{*2}.
\]
On trouve donc
\note{la phrase se poursuit au-delà du feuillet 39, dans le lot suivant}

\end{document}
